\chapter{Experimental setup and methodology}
\label{ch:setup}

Chapters~\ref{ch:background} and~\ref{ch:methods} presented five methods that
approximate the same predictive integral and differ only in the way the posterior
is approximated. This chapter describes the experimental layer that puts that
statement to the test. The scope is deliberately narrow. The experiments are a
demonstration of the compared methods rather than a benchmark study, so the
protocol was designed for interpretability rather than for coverage.

A comparison of posterior approximations is only meaningful if the posterior
approximation is the only quantity that varies. The datasets, the splits, the
preprocessing, the network architecture, the noise model, the prior and the
optimisation settings are therefore identical for every method, and are described
in the sections that follow. The Gaussian process is the one exception, since its
inference is exact and it has no network architecture at all; it is treated as a
reference point rather than as an equal participant.

A deterministic network trained by MAP estimation is included alongside the five
methods. It is not a sixth method but a lower reference for the epistemic term,
which is exactly zero by construction. Any epistemic uncertainty reported by the
remaining methods is thus read against a model that reports none.

% =====================================================================
\section{Experimental setup}
\label{sec:setup-design}

\subsection{Base network and training protocol}
\label{subsec:base-network}

All network-based methods share a single backbone: one hidden layer of 50 units
with a TanH activation, a scalar mean output and one global log-variance parameter.
The width and depth were not chosen freely. A network with one hidden layer of 50
units is the architecture of the protocol of Hernández-Lobato and
Adams~\cite{hernandez-lobato2015}, adopted by the works whose published results are
used for comparison in Section~\ref{sec:results-benchmarks}.

The noise model is homoscedastic for every method. A single variance parameter
describes the whole input space, rather than an input-dependent variance produced
by a second output head. This choice is not driven only by the Laplace implementation
used here, which supports scalar output noise alone. With $\sigma$ independent of $x$,
the aleatoric term cannot widen in sparsely populated regions, so it cannot absorb part
of the epistemic term. The decomposition reported in Chapter~\ref{ch:results} is
therefore protected by the construction of the model and not only by the estimator.
The log-variance is clamped to $[-12, 6]$ for numerical stability, a range wide
enough to be inactive on all datasets used. On the benchmark datasets the variance
parameter is fitted, since the noise level there is one of the quantities every method
has to estimate; on the synthetic problem, where the true value is known by
construction, it is fixed to that value for the network-based methods.

The prior is $\mathcal{N}(0, \gamma^2 I)$ with $\gamma = 1$ for every method. It is
applied as an explicit penalty $\lVert \theta \rVert^2 / (2 \gamma^2 N)$ added to the
loss, with the weight decay of the optimiser set to zero. The alternative, relying on
the \texttt{weight\_decay} argument, would make the effective strength of the prior
depend on the adaptive scaling of Adam and hence on the gradient history of each
parameter. Because the libraries used express the same prior through different
arguments --- a standard deviation, a precision or a penalty coefficient --- the
conversion is implemented once and reused, so that the three parametrisations cannot
drift apart. It should be noted that two methods escape this arrangement for
structural reasons: the implied prior of MC dropout is Bernoulli, and the prior of a
Gaussian process is placed over functions rather than over weights.

Training uses Adam with a learning rate of $10^{-2}$ and a batch size of 128, shared
by all network-based methods. The batch size is a controlled factor rather than a
cost setting, because it fixes the number of optimiser steps per epoch and therefore
the optimisation regime. The number of epochs was determined per dataset, as the
smallest value from a fixed grid at which the mean validation log-likelihood of the
slowest-converging method stopped improving. Taking the maximum over methods
guarantees that no method is under-trained on any dataset.
Table~\ref{tab:training-protocol} collects the resulting settings.

\begin{table}[htbp]
    \centering
    \caption{Shared training protocol.}
    \label{tab:training-protocol}
    \begin{tabular}{ll}
        \hline
        Setting & Value \\
        \hline
        Hidden layers $\times$ units & $1 \times 50$ \\
        Activation                   & TanH \\
        Output                       & mean, plus one global $\log \sigma^2$ \\
        Prior                        & $\mathcal{N}(0, I)$, explicit loss penalty \\
        Optimiser                    & Adam, learning rate $10^{-2}$ \\
        Batch size                   & 128 \\
        Epochs                       & per dataset, see Table~\ref{tab:epochs} \\
        \hline
    \end{tabular}
\end{table}

\begin{table}[htbp]
    \centering
    \caption{Number of training epochs used for each dataset.}
    \label{tab:epochs}
    \begin{tabular}{lc}
        \hline
        Dataset & Epochs \\
        \hline
        \texttt{sin\_homo}         & 2000 \\
        \texttt{yacht}             & 2000 \\
        \texttt{energy}            & 1000 \\
        \texttt{concrete}          & 1000 \\
        \texttt{wine\_quality\_red}& 20   \\
        \texttt{kin8nm}            & 1000 \\
        \texttt{power\_plant}      & 50   \\
        \hline
    \end{tabular}
\end{table}

One limitation of this rule is worth stating. On \texttt{yacht} the budget was
capped at the value given above, because the two methods that set it move in
opposite directions: the deterministic baseline reaches its optimum early and
degrades afterwards, while MC dropout keeps improving to the end of every grid
that was measured. Extending the grid would therefore raise the shared budget
indefinitely and penalise the baseline further, so the optimum of MC dropout on
that dataset remains unmeasured.

\subsection{Method configurations}
\label{subsec:method-configurations}

Table~\ref{tab:method-config} lists the configuration of each method. Only the
settings that are specific to a method appear there; everything shared is given in
Table~\ref{tab:training-protocol}. Established implementations were preferred to
custom code wherever one exists, so that each method reproduces the algorithm as
published rather than a local variant of it.

\begin{table}[htbp]
    \centering
    \caption{Per-method configuration. Shared settings are omitted.}
    \label{tab:method-config}
    \begin{tabular}{lll}
        \hline
        Method & Implementation & Settings \\
        \hline
        MAP baseline       & custom            & deterministic, zero epistemic term \\
        Gaussian process   & \texttt{scikit-learn} & $C \times \mathrm{RBF} + \mathrm{White}$, 5 restarts \\
        Bayes by Backprop  & \texttt{bayesian-torch} & mean-field Gaussian, $T = 100$ \\
        MC dropout         & custom            & $p = 0.1$, $T = 100$ \\
        Laplace            & \texttt{laplace-torch}  & full Hessian, linearised prediction \\
        Deep ensembles     & custom            & $M = 5$ members \\
        \hline
    \end{tabular}
\end{table}

The Gaussian process uses a constant kernel multiplied by a radial basis function
kernel, with an additive white-noise term. The epistemic variance is the posterior
variance of the latent function and the aleatoric variance is the fitted noise level,
which is constant across the input space. Target standardisation is performed outside
the model, identically to the networks, so that the reported log-likelihoods remain on
the same scale.

Bayes by Backprop is realised with reparameterised linear layers, which implement the
algorithm of Blundell et al.~\cite{blundell2015} directly. The Kullback--Leibler term
is weighted uniformly across minibatches, and the predictive moments are estimated from
$T$ stochastic passes, with the divisor $T - 1$ in the epistemic term as in
Equation~\eqref{eq:bbb-variance-estimator}. The data term of the objective is averaged
over $32$ reparameterisation samples per step, in every experiment. The value is not a
cost setting but a property of the estimate: with a single sample the variational
standard deviations barely move from their initialisation, and raising the count to $32$
increases the median epistemic standard deviation by a factor of roughly seven on the
synthetic problem.

MC dropout keeps the masks active at prediction time, with the dropout rate fixed at
$p = 0.1$ and the number of passes at $T = 100$. Both values are held constant across
all experiments, since they govern the estimate itself rather than only its cost.
The masks are applied to the hidden activations alone, which is a deliberate departure
from the derivation of Section~\ref{sec:mc-dropout}, where dropout precedes every
weighted layer. On a one-dimensional input, a mask on the input layer removes the only
feature the network has in a fraction of the forward passes, at training and at
prediction time alike. The network compensates by inflating its noise estimate to cover
the resulting loss spikes, so the aleatoric term absorbs an artefact of the masking
rather than the noise of the data. Measured on the synthetic problem with the noise
parameter left free, where the true variance is $0.01$, the mean aleatoric variance was
$0.103$ with input dropout and $0.039$ without it. The literature-faithful variant is
restored in the separate run that validates the implementation against published
results, so that this departure cannot serve as an alternative explanation of any
discrepancy found there.

The Laplace approximation is fitted over all network weights, with a full covariance
matrix and a linearised predictive distribution, following
Equation~\eqref{eq:laplace-predictive}. The full structure is affordable here because
the network has a few hundred parameters. The prior precision is fixed at $1/\gamma^2$
rather than tuned by maximising the marginal likelihood; tuning it would let this one
method adopt a prior different from the shared one, which is precisely the asymmetry
the setup is meant to avoid.

Deep ensembles consist of $M = 5$ networks that differ in their initialisation and in
the order of their minibatches. The method therefore receives $M$ times the training
budget of every other one. Equalising the budget would require narrowing the individual
members and would break the agreement with the reference architecture, so the choice made
here is one of equal architecture rather than equal computation, and the training time is
reported as a separate metric. Adversarial training was not applied. Lakshminarayanan
et al.~\cite{lakshminarayanan2017} report it as an optional refinement, and including
it would add to this method a component absent from the other four. The predictive
moments are computed with the same estimator as for Bayes by Backprop, with $T = M$.

% =====================================================================
\section{Datasets}
\label{sec:datasets}

\subsection{Synthetic problem}
\label{subsec:synthetic-data}

The synthetic problem is one-dimensional and its data-generating process is known,
which makes it the only setting in which the estimated uncertainty can be compared
with the truth. The target is $y = \sin(x) + \varepsilon$ with
$\varepsilon \sim \mathcal{N}(0, 0.1^2)$. Training points are placed on a regular grid
over $[0, 6]$ and predictions are evaluated on $[-2, 8]$, so that a substantial part of
the evaluation range lies outside the training range. Extrapolation is thus enforced
rather than hoped for, and the epistemic term has a region in which it is expected to
grow. The evaluation range is symmetric about the training interval, extending two
units beyond it on each side. An asymmetric range would let a method appear better on
the side where it is simply closer to the data.

The training set contains 250 points. A smaller set would be closer to the usual
illustrations, but it was found to be unusable here. The network has of the order of
one hundred and fifty parameters, so with a few dozen observations a large part of
the parameter space receives no information from the data and remains constrained by
the prior alone. The result is narrow spikes in the epistemic variance of the Laplace
approximation inside the training range, which are an artefact of the sample size
rather than a property of the method. With 250 points the effect disappears.

\subsection{Benchmark datasets}
\label{subsec:benchmark-data}

Six regression datasets from the UCI repository were used, listed in
Table~\ref{tab:datasets}. They occupy the smaller end of the collection commonly
reported in this literature, which is a consequence of a deliberate choice: the sets
were selected so that an exact Gaussian process remains feasible on the full training
fold. The alternative, subsampling the data for the Gaussian process alone, would have
placed one method on a different training set from the others.

\begin{table}[htbp]
    \centering
    \caption{Benchmark datasets. $N$ is the total number of observations, $N_{\mathrm{train}}$
    the size of the training fold and $d$ the input dimensionality.}
    \label{tab:datasets}
    \begin{tabular}{lccc}
        \hline
        Dataset & $N$ & $N_{\mathrm{train}}$ & $d$ \\
        \hline
        \texttt{yacht}              & 308  & 277  & 6  \\
        \texttt{energy}             & 768  & 691  & 8  \\
        \texttt{concrete}           & 1030 & 927  & 8  \\
        \texttt{wine\_quality\_red} & 1599 & 1439 & 11 \\
        \texttt{kin8nm}             & 8192 & 7373 & 8  \\
        \texttt{power\_plant}       & 9568 & 8611 & 4  \\
        \hline
    \end{tabular}
\end{table}

Each dataset is evaluated on the twenty random $90/10$ splits distributed with the
reference implementation of MC dropout~\cite{gal2016}, rather than on splits generated
here. The datasets are small enough that independently generated splits produce results
that cannot be placed next to published ones, and the comparison in
Section~\ref{sec:results-benchmarks} depends on that being possible. Using the
distributed index files also settles several choices that would otherwise be arbitrary,
among them which target variable is used for \texttt{energy} and which variant of the
wine dataset is meant.

The \texttt{boston} housing dataset appears in most published tables and was
nevertheless omitted. It has been withdrawn from \texttt{scikit-learn} because one of
its features encodes a racial characteristic. The consequence is stated plainly: one
row of reference values fewer is available than in the works cited.

Both the inputs and the target are standardised, with the transformation fitted on the
training fold alone. All metrics are computed after the transformation of the target
has been inverted, and the inversion is applied to the whole predictive distribution
rather than to its mean, so that the reported standard deviations are also expressed in
the original units.

Exact Gaussian process inference was performed on the four smaller datasets. On
\texttt{kin8nm} and \texttt{power\_plant} it is not reported. The method remains
feasible at those sizes in principle, but the product of twenty splits and repeated
restarts of the kernel optimiser raises the cost out of proportion to what the
resulting row would add. The decision rests on a measurement rather than on an
assertion: a single split of \texttt{kin8nm} was fitted as a cost probe and took
$31$ minutes against a limit of twenty. The corresponding cells are marked as not run, and every
comparative statement about the Gaussian process in Chapter~\ref{ch:results} is
restricted to the range of $N$ in which it was actually evaluated.

% =====================================================================
\section{Evaluation metrics}
\label{sec:metrics}

Point accuracy is measured by the root mean squared error over the $n$ test points,
%
\begin{equation}
    \mathrm{RMSE} = \sqrt{\frac{1}{n} \sum_{i=1}^{n} \bigl(y_i - \hat{\mu}(x_i)\bigr)^2 } ,
    \label{eq:rmse}
\end{equation}
%
which depends on the predictive mean alone and says nothing about the uncertainty. The
quality of the full predictive distribution is measured by the average test
log-likelihood,
%
\begin{equation}
    \mathrm{LL} = \frac{1}{n} \sum_{i=1}^{n}
      \log \mathcal{N}\!\left(y_i \mid \hat{\mu}(x_i), \hat{\sigma}^2(x_i)\right) ,
    \label{eq:ll}
\end{equation}
%
reported in the convention of the literature, in which larger values are better and the
values themselves are negative. A model whose predictive variance is too small is
penalised here even when its mean is accurate, which is the property that makes this
the central metric of the comparison.

Two interval metrics are reported alongside. The prediction interval coverage
probability at the nominal level of $95\%$,
%
\begin{equation}
    \mathrm{PICP} = \frac{1}{n} \sum_{i=1}^{n}
      \mathbf{1}\!\left[ l_i \le y_i \le u_i \right] ,
    \label{eq:picp}
\end{equation}
%
gives the fraction of test targets that fall inside the interval $[l_i, u_i]$, and the
mean prediction interval width,
%
\begin{equation}
    \mathrm{MPIW} = \frac{1}{n} \sum_{i=1}^{n} (u_i - l_i) ,
    \label{eq:mpiw}
\end{equation}
%
gives its average width. Neither is informative on its own. Coverage can be raised to
any level by widening the intervals, and narrow intervals are only a merit if the
coverage is retained, so the two are always reported together.

Training time is reported as a fourth quantity. It is not a measure of quality, but it
is required by the budget asymmetry described in
Section~\ref{subsec:method-configurations}: deep ensembles are trained five times over, and
a comparison that omitted this would present the cost as free. The reported times come
from a separate measurement run, executed sequentially on one split with the number of
threads fixed. The times recorded during the main run are not used, because that run was
executed in parallel and the resulting figures are inflated by the shared load.

Of these quantities, RMSE and the log-likelihood are the two that appear in published
tables and are therefore directly comparable with them. The remaining metrics extend the
comparison but have no reference values. This division is used in
Chapter~\ref{ch:results}: the first two columns serve to validate the implementation
against published results for MC dropout, and the others describe behaviour for which
this work provides the only measurements.
